%%%%%%%%%%%%%%%%%%%%%%%%%%%%%%%%%%%%%%%%%%%%%%%%%%%%%%%%%%%%%%%%%%%%%%%%%%%
% Handout und Uebungsblatt zu Lehreinheit 2
% Im Skript am Kapitelende, in handouts.tex als Druckvorlage.
%%%%%%%%%%%%%%%%%%%%%%%%%%%%%%%%%%%%%%%%%%%%%%%%%%%%%%%%%%%%%%%%%%%%%%%%%%%
\begin{handout}{Die Forschungsfrage}

\textbf{Worum es heute ging.} \enquote{Unter diesem Becher liegt etwas, das
noch kein Mensch je gesehen hat.} Der Satz stimmt: Es ist eine Erdnuss in
der Schale. Wir haben ihn aber als Aussage über eine unbekannte \emph{Art}
von Ding gelesen. Jede Frage bringt solche stillen Voraussetzungen mit. Eine
Forschungsfrage legt sie offen und sagt, welche Beobachtung sie beantworten
würde.

\textbf{Kernbegriffe}
\begin{description}
  \item[Forschungsfrage] legt fest, welcher Sachverhalt bei welchem
    Gegenstand unter welchen Bedingungen geklärt werden soll, so genau, dass
    man weiß, welche Daten sie beantworten.
  \item[Informationslücke] Neugier entsteht, wenn man eine begrenzte Lücke
    zwischen Wissen und Wissenwollen bemerkt \parencite{loewenstein1994}.
  \item[Implizite Prämisse] Annahme, die eine Frage voraussetzt, ohne sie
    auszusprechen (\enquote{Warum schadet X?} setzt einen Schaden voraus).
  \item[Operationalisierung] Angabe, durch welche beobachtbaren Größen und
    nach welcher Regel ein Begriff gemessen wird
    \parencite{doering2016forschungsmethoden}.
  \item[Forschungslücke] Teil des Gegenstands, zu dem es noch keine
    belastbare Antwort gibt; man findet sie durch Recherche.
\end{description}

\textbf{Vom Thema zur Frage: die Probe nach \textcite{booth2016craft}}\\
Ich untersuche \emph{X}, weil ich herausfinden will, \emph{Y}, damit andere
\emph{Z} verstehen.

\textbf{SMART-Prüfung einer Forschungsfrage}

\begin{center}
\small
\begin{tabular}{lp{0.7\textwidth}}
\toprule
\textbf{S}pezifisch & ein Gegenstand, eine Population, ein Kontext \\
\textbf{M}essbar & klar, welche Größe erhoben wird \\
\textbf{A}chievable (erreichbar) & Zugang zu Daten, Personen, Material \\
\textbf{R}elevant & die Antwort verändert Wissen oder Praxis \\
\textbf{T}erminiert & in der verfügbaren Zeit zu schaffen \\
\bottomrule
\end{tabular}
\end{center}

\textbf{Drei Prüffragen gegen versteckte Prämissen:} (1) Was muss wahr sein,
damit die Frage Sinn ergibt? (2) Welche Antwort schließt die Frage aus?
(3) Welches Wort würde jemand anderes anders verstehen?

\begin{merkbox}
Eine Forschungsfrage ist eine Informationslücke mit Rand. Wer die Prämissen
seiner Frage nicht kennt, findet am Ende die Antwort, die er hineingelegt
hat.
\end{merkbox}

\textbf{Zum Weiterlesen:} \textcite{doering2016forschungsmethoden},
Kapitel zu Forschungsproblem und Operationalisierung;
\textcite{booth2016craft}, Teil~I.
\end{handout}

\begin{uebungsblatt}{Fragen sezieren}

\aufgabe{Verborgene Prämissen finden}[in der Stunde, Zweiergruppe, 10 Min.]
Die folgenden Fragen sind konstruiert. Jede enthält mindestens eine
Voraussetzung, die nicht ausgesprochen wird. Wählen Sie zwei Fragen, notieren
Sie die Prämisse(n) und formulieren Sie eine Fassung, die auch eine
unerwartete Antwort zulässt.

\begin{enumerate}
  \item Warum lernen Studierende mit Laptop in der Vorlesung weniger als mit
    Papier und Stift?
  \item Wie stark beeinflusst der Vollmond die Zahl der Notaufnahmen?
  \item Welche Strategien nutzen erfolgreiche Gründer:innen, um Rückschläge
    zu verkraften?
  \item Ab welchem Alter verlieren Kinder ihre Kreativität?
  \item Wie viel Prozent ihres Gehirns nutzen Menschen im Alltag?
  \item Warum sind Frauen die besseren Führungskräfte?
  \item Wie kann die Hochschule verhindern, dass ChatGPT das eigenständige
    Denken zerstört?
  \item Wie viele Stunden Schlaf braucht ein Mensch, um leistungsfähig zu
    sein?
\end{enumerate}

\begin{center}
\renewcommand{\arraystretch}{2.2}
\begin{tabular}{|p{0.06\textwidth}|p{0.38\textwidth}|p{0.46\textwidth}|}
\hline
\textbf{Nr.} & \textbf{verborgene Prämisse(n)} & \textbf{offenere Fassung} \\
\hline
 & & \\ \hline
 & & \\ \hline
\end{tabular}
\end{center}

\aufgabe{Die eigene Frage sezieren}[in der Stunde, Zweiergruppe, 20 Min.]
Schreiben Sie Ihre mitgebrachte Frage auf. Ihr Gegenüber stellt die drei
Prüffragen; Sie notieren die Antworten. Wer keine eigene Frage hat, wählt
ein Ersatzthema: Mensaessen, Pendeln zur Hochschule, Lernen mit Karteikarten,
Nebenjobs im Studium.

\smallskip
\noindent\textbf{Meine Frage (Rohfassung):}
\schreiblinien{2}

\begin{center}
\renewcommand{\arraystretch}{2.0}
\begin{tabular}{|p{0.36\textwidth}|p{0.56\textwidth}|}
\hline
Was wird vorausgesetzt? & \\ \hline
Welche Wörter sind unscharf? & \\ \hline
Was würde ich messen (Indikator)? & \\ \hline
Bei wem, wo, wann? & \\ \hline
Welche Antwort würde mich überraschen? & \\ \hline
\end{tabular}
\end{center}

\aufgabe{SMART-Check}[in der Stunde, Zweiergruppe]
Formulieren Sie die Frage neu und kreuzen Sie an, welche Kriterien sie
erfüllt.
\schreiblinien{2}
\noindent$\square$ spezifisch \quad $\square$ messbar \quad
$\square$ erreichbar \quad $\square$ relevant \quad $\square$ terminiert

\aufgabe{Forschungsfrage schriftlich}[Hausaufgabe, Abgabe bis LE~3]
Formulieren Sie eine wissenschaftliche Forschungsfrage zu Ihrem
Semesterthema. Begründen Sie in drei Sätzen, warum diese Frage beantwortbar
und nicht trivial ist. Füllen Sie zusätzlich die Probe aus: \enquote{Ich
untersuche \dots, weil ich herausfinden will, \dots, damit andere \dots\
verstehen.}
\schreiblinien{5}

\aufgabe{Tagebuch des Staunens}[Hausaufgabe]
Formulieren Sie Ihr Staunerlebnis dieser Woche als Frage und notieren Sie
darunter die Prämisse, die in ihr steckt.

\begin{lehrhinweis}[Erwartungshorizont]
Aufgabe~1, mögliche Prämissen: (1) dass Laptop-Nutzer weniger lernen
(Richtung vorausgesetzt) und dass \enquote{Lernen} klar gemessen ist;
(2) dass es einen Einfluss gibt; (3) dass \enquote{erfolgreich} definiert ist
und Erfolg mit den Strategien zusammenhängt (Überlebenden-Auswahl: die
Gescheiterten fehlen); (4) dass Kinder Kreativität verlieren und dass sie
sich als eine Größe messen lässt; (5) setzt einen festen Prozentsatz voraus,
die bekannte Zehn-Prozent-Legende; (6) Richtung und Gruppenvergleich
vorausgesetzt, \enquote{besser} undefiniert; (7) dass das Denken zerstört
wird und dass die Hochschule das verhindern kann; (8) dass es eine Zahl für
alle Menschen gibt und \enquote{leistungsfähig} klar ist. Offenere Fassungen
fragen nach \emph{ob} und \emph{wie stark} statt nach \emph{warum}, nennen
Population und Messgröße. Aufgabe~2/3: Die Frage muss am Ende Population,
Messgröße und Zeitraum enthalten; eine Warum-Frage darf bleiben, wenn die
Prämisse vorher geprüft wird. Aufgabe~4: Begründung nennt Datenquelle und
zeigt, dass die Antwort nicht offensichtlich ist.
\end{lehrhinweis}
\end{uebungsblatt}
